\newpage
\section{Introduction}
\label{sec:introduction}

\subsection{Overview}
\label{subsec:intro-overview}

6G networking is expected to support ultra-low latency services, massive connectivity, and deeply integrated intelligent control across communication and computing infrastructure. In such systems, computation can no longer be treated as a centralized cloud-only service. Instead, latency-sensitive and energy-constrained workloads must be distributed over a hierarchical computing continuum in which local devices, edge data centers, and remote cloud servers cooperate in real time.

This report studies task offloading and wireless PRB allocation under a three-tier LOCAL--EDGE--CLOUD architecture. The core difficulty is that offloading decisions are coupled with communication resource allocation. Selecting a distant yet powerful cloud server may reduce execution delay but can increase transmission time and bandwidth contention, while aggressive local execution preserves network resources but may overload weak devices. The resulting optimization problem is high-dimensional, non-stationary, and strongly dependent on interactions among multiple devices.

To address this challenge, the project formulates the problem as a cooperative multi-agent decision process and proposes an on-policy Multi-Agent Proximal Policy Optimization framework. The main contributions are summarized as follows.
\begin{enumerate}[(a)]
    \item The joint offloading and PRB allocation problem is formulated as a Dec-POMDP in which all task-generating devices act as decentralized agents.
    \item A MAPPO framework is designed with per-device actors augmented by learnable agent embeddings and a centralized critic operating on a 27-dimensional global state.
    \item A dual-head discrete action space is introduced to jointly decide the offloading destination and the PRB level assigned to each task.
    \item PureEdgeSim is extended into a complete training, inference, and evaluation pipeline with PRB-aware networking, a distance-sensitive channel model, TCP JSON co-simulation, telemetry logging, and comparison against heuristic and RL baselines.
\end{enumerate}

\subsection{Background}
\label{subsec:intro-background}

\subsubsection{6G and Mobile Edge Computing}

The convergence of communication and computing is one of the defining characteristics of 6G visions and beyond-5G architectures. Edge computing pushes computation closer to users and sensors, reducing end-to-end latency and backhaul pressure while enabling context-aware services such as augmented reality, connected healthcare, and industrial automation \cite{satyanarayanan2017emergence, mao2017survey, saad2020vision}. In practice, however, the edge alone is insufficient for all workloads. A realistic deployment requires a hierarchical system in which local devices process lightweight or privacy-sensitive tasks, nearby edge data centers handle delay-critical workloads, and cloud servers provide elastic capacity for compute-intensive requests.

\subsubsection{Traditional Task Offloading Methods}

Classical offloading strategies typically rely on rule-based heuristics, convex optimization, or game-theoretic formulations. Greedy and round-robin schemes are simple to implement but ignore long-term system balance. Optimization-based methods can achieve strong performance under static assumptions, yet they often require accurate models, centralized observability, or repeated re-solving when the environment changes. These assumptions become restrictive in mobile multi-user settings with heterogeneous devices and fluctuating wireless conditions.

\subsubsection{Why Multi-Agent Reinforcement Learning}

Reinforcement learning offers an attractive alternative because it can adapt online through repeated interaction with the environment. Nevertheless, single-agent formulations suffer from state-space explosion when many devices compete for shared edge and wireless resources. They also fail to model device-specific behaviors and coordination patterns explicitly. MARL mitigates these issues by assigning local policies to individual devices while preserving coordinated learning through centralized critics or shared training signals \cite{lowe2017maddpg, yu2022surprising}.

\subsection{Related Work}
\label{subsec:intro-related}

Recent studies have applied deep reinforcement learning to mobile edge computing for task offloading, server selection, and radio resource management. Many early approaches focused on single-agent DQN or PPO policies and treated the environment as centrally observable. While effective in small settings, they become increasingly brittle as the number of devices and feasible destinations grows. More recent work explores cooperative MARL, including MADDPG, QMIX, and MAPPO-style centralized-training approaches, to capture interactions among users and edge nodes.

In parallel, simulation platforms have become a critical component of research on intelligent edge orchestration. PureEdgeSim provides a lightweight discrete-event simulator for cloud, edge, and mist computing \cite{mechalikh2021pureedgesim}. CloudSim Plus serves as its underlying cloud simulation engine \cite{silvafilho2017cloudsimplus}. EISim extends PureEdgeSim toward DRL-enabled experimentation \cite{abid2023eisim}, while the ML-RL-PureEdgeSim codebase adds an initial learning-oriented structure around the original simulator \cite{mlrlpureedgesimgithub}. The present project builds on that direction and introduces a richer MARL environment interface, PRB-aware networking, a stronger actor design, and a bidirectional Java-Python runtime.

\begin{table}[htbp]
    \centering
    \begin{tabular}{p{0.20\linewidth}p{0.18\linewidth}p{0.14\linewidth}p{0.14\linewidth}p{0.24\linewidth}}
        \hline
        \textbf{Work} & \textbf{Simulator} & \textbf{Joint Offloading + Radio} & \textbf{MARL} & \textbf{Key Limitation or Difference} \\
        \hline
        Representative DRL offloading studies & custom / abstract models & partial & limited & often assume centralized state or simplified links \\
        EISim-style DRL simulation & PureEdgeSim-derived & partial & optional & emphasizes DRL integration but not the present PRB design \\
        ML-RL-PureEdgeSim baseline & PureEdgeSim & limited & initial RL support & does not provide the full MAPPO device-agent pipeline used here \\
        This work & PureEdgeSim + custom extensions & yes & yes & joint destination and PRB allocation with CTDE and co-simulation \\
        \hline
    \end{tabular}
    \caption{Positioning of this report with respect to related edge-computing simulation and RL work.}
    \label{tab:related-work-comparison}
\end{table}

\subsection{Report Organization}
\label{subsec:intro-organization}

The remainder of this report is organized as follows. \Secref{sec:system-model} introduces the network, task, communication, computation, and energy models together with the Dec-POMDP formulation. \Secref{sec:simulation} describes the customized PureEdgeSim environment and the Java-Python co-simulation architecture. \Secref{sec:methodology} details the proposed MAPPO framework, observation and action design, network architectures, reward function, and training algorithm. \Secref{sec:results} presents the experimental setup, training behavior, multi-seed evaluation, and qualitative comparison with representative baselines. Finally, \Secref{sec:conclusion} concludes the report and outlines future work.
